%! Tex main = MarquezSalazarBrandon-PropuestaDeTesis.tex
\subsection{Antecedentes}

\subsubsection{NIM en la investigación}

El ámbito de la investigación de alto rigor, como el que
caracteriza a la química o la cosmología, suele asociarse con
métodos matemáticos formalmente fundamentados. Sin embargo, en
las últimas décadas, los métodos metaheurísticos han irrumpido
en estos campos con resultados notables. Su aplicación has
permitido demostrar que, a pesar de carecer de la formalidad
matemática teórica que tradicionalmente se exige en estas
disciplinas, resultan ser herramientas perfectamente aplicables
y, en muchos contextos, altamente efectivas \cite{Wang_2010, Prasad_2014}.

Esta efectividad se traduce en logros concretos: una reducción
sustancial del coste computacional y, cuando el modelo está
bien sintonizado y la función objetivo correctamente definida,
una precisión superior gracias a sus capacidades de explotación,
llegando incluso a superar los valores obtenidos mediante
enfoques tradicionales \cite{Megantoro_2025}.

\subsubsection{Aplicaciones metaheurísticas en la exploración de configuraciones de vacío}

El uso de metaheurísticas para explorar el espacio de configuraciones en
teoría de cuerdas ha demostrado mejoras significativas en eficiencia,
capacidad de descubrimiento y precisión respecto a los métodos de búsqueda
aleatoria. A continuación se reseñan los principales avances agrupados por
las cualidades que dichas técnicas han evidenciado.


\paragraph{Eficiencia en la búsqueda}

El trabajo pionero de \textcite{Abel_2014} estableció que los algoritmos
genéticos son órdenes de magnitud más eficientes que el escaneo aleatorio
para localizar configuraciones con propiedades fenomenológicas específicas
dentro de modelos fermiónicos libres. Este resultado fue relevante para
demostrar que, aunque el espacio de búsqueda es inmenso (del orden de
\(2^{51}\) configuraciones posibles), la estructura del paisaje permite
que métodos guiados por selección y cruce converjan rápidamente hacia
regiones de interés.

Posteriormente, el trabajo de \textcite{Bizet_2020} reportó más de sesenta mil
configuraciones de flujo con puntos críticos en compactificaciones Tipo
IIB, empleando un algoritmo genético acoplado a una red neuronal. Este
enfoque híbrido permitió reducir drásticamente el tiempo de cómputo
respecto a la fuerza bruta, evidenciando que la sinergia entre
metaheurísticas y aprendizaje automático puede escalar la exploración
a espacios de mayor complejidad.

\paragraph{Capacidad de descubrimiento y generación de nuevas configuraciones}

Más allá de la búsqueda eficiente, las metaheurísticas han demostrado
capacidad para generar configuraciones novedosas no reportadas
previamente. \cite{He_2025} utilizó templado simulado
(\foreignlanguage{english}{\textit{simulated annealing}}) para construir
\foreignlanguage{english}{\textit{brane tilings}} geométricamente
consistentes, logrando descubrir un nuevo
\foreignlanguage{english}{\textit{tiling}} con 26 campos cuánticos que no
aparecía en catálogos exhaustivos previos. Este trabajo evidencia que las
metaheurísticas no solo optimizan dentro de espacios conocidos, sino que
pueden explorar regiones inéditas del espacio de soluciones.

\paragraph{Precisión y calidad de las soluciones}
 
En el ámbito de la clasificación geométrica, \cite{Gao_2022} alcanzó
precisiones superiores al 99.9\% en la predicción de politopos que
conducen a orientifolds de Calabi-Yau y vacíos Tipo IIB. Aunque su enfoque
principal es una red neuronal convolucional, el proceso de entrenamiento
y validación cruzada demuestra que los métodos computacionales pueden
capturar patrones sutiles en los datos de entrada (las coordenadas de los
politopos) y traducirlos en predicciones de alta calidad sobre la
viabilidad de las configuraciones resultantes.

\paragraph{Exploración de regiones de transición}

Un avance particularmente relevante para el problema de optimización de
potenciales es el reportado por \cite{Lanza_2025}. Mediante técnicas de
aprendizaje automático, los autores identificaron regiones de "penumbra"
(zonas de transición entre el interior del espacio de módulos y sus
límites asintóticos) donde aparecen configuraciones metaestables con
energía positiva. Aunque su enfoque no es puramente metaheurístico, el
trabajo demuestra que dichas regiones albergan estructuras finas del
potencial que son susceptibles de ser exploradas mediante métodos de
optimización global.

\paragraph{Contexto metodológico complementario}

Finalmente, \cite{De_Luca_2025} abordó el problema desde una perspectiva
numérica diferente: la resolución directa de ecuaciones diferenciales
parciales mediante redes neuronales. Si bien no emplea metaheurísticas, su
trabajo contextualiza la necesidad de métodos computacionales robustos
para atacar problemas geométricos en compactificaciones, y establece un
punto de comparación respecto a enfoques alternativos para construir
métricas y configuraciones de vacío.

\paragraph{Hacia una extensión del enfoque}

Los trabajos aquí reseñados muestran una tendencia clara: el uso de
metaheurísticas permite explorar eficientemente el paisaje de vacíos. En
particular, \cite{Damian_2022} implementó templado simulado para
identificar más de 170 configuraciones dS estables en compactificaciones
Tipo IIB con torsión, demostrando la viabilidad del método. No obstante,
dicho estudio se limitó a un único algoritmo estocástico. El presente
trabajo se propone extender este análisis mediante la implementación
y comparación de un conjunto más amplio de metaheurísticas, con el fin de
evaluar su desempeño relativo y determinar si es posible incrementar la
cantidad de vacíos estables identificados respecto a lo reportado en
\cite{Damian_2022}.
